\documentclass[10pt]{article}
% Shared typographic and algebraic source.
\usepackage[letterpaper,margin=0.86in]{geometry}
\usepackage[T1]{fontenc}
\usepackage[utf8]{inputenc}
\usepackage{lmodern}
\usepackage{microtype}
\usepackage{amsmath,amssymb,amsthm,mathtools,mathrsfs}
\usepackage{booktabs,array,tabularx,longtable}
\usepackage{enumitem}
\usepackage[dvipsnames]{xcolor}
\usepackage{fancyhdr,titlesec}
\usepackage[hidelinks]{hyperref}

\definecolor{ink}{HTML}{1F2A33}
\definecolor{teal}{HTML}{226A73}
\definecolor{muted}{HTML}{65747E}
\definecolor{line}{HTML}{C5D1D3}
\definecolor{wash}{HTML}{F3F7F6}

\setlength{\parindent}{0pt}
\setlength{\parskip}{0.46em}
\setlength{\emergencystretch}{2em}
\allowdisplaybreaks[2]
\renewcommand{\arraystretch}{1.16}
\setlist{nosep,leftmargin=1.5em}

\pagestyle{fancy}
\fancyhf{}
\fancyhead[L]{\footnotesize\sffamily\color{muted}COPPALINI ARCHITECTURE}
\fancyhead[R]{\footnotesize\sffamily\color{muted}PHYSICAL REALITY INSIDE CONSTRAINT}
\fancyfoot[C]{\footnotesize\sffamily\color{muted}\thepage}
\renewcommand{\headrulewidth}{0.35pt}
\renewcommand{\headrule}{\hbox to\headwidth{\color{line}\leaders\hrule height \headrulewidth\hfill}}
\setlength{\headheight}{19pt}

\titleformat{\section}{\Large\bfseries\color{ink}}{\thesection.}{0.55em}{}
\titleformat{\subsection}{\large\bfseries\color{teal}}{\thesubsection.}{0.5em}{}
\titlespacing*{\section}{0pt}{1.42em}{0.48em}
\titlespacing*{\subsection}{0pt}{1.02em}{0.32em}

\newtheorem{theorem}{Theorem}
\newtheorem{proposition}[theorem]{Proposition}
\newtheorem{lemma}[theorem]{Lemma}
\theoremstyle{definition}
\newtheorem{definition}[theorem]{Definition}
\theoremstyle{remark}
\newtheorem{remark}[theorem]{Boundary}

\newcommand{\RR}{\mathbb R}
\newcommand{\CC}{\mathbb C}
\newcommand{\QQ}{\mathbb Q}
\newcommand{\ZZ}{\mathbb Z}
\newcommand{\Id}{I}
\newcommand{\tr}{\operatorname{tr}}
\newcommand{\im}{\operatorname{im}}
\newcommand{\kerop}{\operatorname{ker}}
\newcommand{\Span}{\operatorname{span}}
\newcommand{\diag}{\operatorname{diag}}
\newcommand{\snoc}{\operatorname{snoc}}
\newcommand{\Admit}{\operatorname{Admit}}
\newcommand{\Gate}{\operatorname{Gate}}
\newcommand{\Typed}{\operatorname{Typed}}
\newcommand{\Future}{\operatorname{future}}
\newcommand{\Scope}{\operatorname{sc}}
\newcommand{\SilScope}{\mathsf{Silence}_{\rm scope}}
\newcommand{\SilFit}{\mathsf{Silence}_{\rm fit}}

\newcommand{\papertitle}[2]{%
  \begin{center}
  {\color{teal}\rule{\textwidth}{1.35pt}}\\[0.38in]
  {\footnotesize\sffamily\bfseries\color{teal} COPPALINI ARCHITECTURE}\\[0.16in]
  {\fontsize{31}{36}\selectfont\bfseries\color{ink}#1\par}
  \vspace{0.15in}
  {\Large\color{teal}#2\par}
  \vspace{0.22in}
  {\color{line}\rule{\textwidth}{0.55pt}}\\[0.12in]
  {\normalsize Anthony Vito Coppa}\par
  {\small 9 November 2026}\par
  \vspace{0.30in}
  \end{center}
}

\newcommand{\sharedroot}{%
Let
\[
Q=\begin{pmatrix}1&1\\1&0\end{pmatrix},\qquad
\Omega=\begin{pmatrix}0&1\\-1&0\end{pmatrix},\qquad
C=-\Omega,
\]
\[
G=[Q,\Omega]=\begin{pmatrix}-2&1\\1&2\end{pmatrix},\qquad
H=\frac{G}{\sqrt5},\qquad
S=\frac{2Q-I}{\sqrt5}.
\]
The common exact identities are
\[
Q^2=Q+I,\qquad C^2=-I,\qquad G^2=5I,\qquad H^2=S^2=I,
\]
\[
H=CS,\qquad SG=\sqrt5\Omega=-GS,\qquad
L_s:=e^{sS},\qquad
L_s^{\mathsf T}GL_s=G,\qquad HL_sH=L_{-s}.
\]
}

\newcommand{\sharedoccurrence}[2]{%
A physical field has an interface
\[
\begin{aligned}
\mathfrak I_J:=(&J,\Gamma_J,\operatorname{Src}_J;
A_J,X_J,Y_J,T_J,\Sigma_J,\operatorname{Con}_J,\\
&\rho_J,O_J,K_J,B_J,\theta_J),
\qquad
\rho_J:X_J\to Y_J,\quad
O_J:X_J\to\mathcal O_J,\quad
\theta_J:X_J\to\{0,1\}.
\end{aligned}
\tag{#1}
\]
Its occurrence type is
\[
\begin{aligned}
\mathsf{Occ}_J:=\{&
(\mathfrak I_J;\Theta,x,\rho_J(x),\kappa,
\boldsymbol\varepsilon,\theta_J(x)):\\
&\Theta\in A_J,\ x\in X_J,\ \kappa\in K_J,\
\boldsymbol\varepsilon\in\RR_{\ge0}^{3}\}.
\end{aligned}
\tag{#2}
\]
An indexed occurrence $\mathsf{Occ}_{J,n}$ is an element of
$\mathsf{Occ}_J$.
}

\newcommand{\sharedmatching}{%
Let
\[
\mathcal A_8:=\{\infty\}\sqcup\ZZ_7,
\qquad
(\infty,0,1,2,3,4,5,6)
=(\mathsf{Gr},\mathsf{Li},\mathsf{Me},\mathsf{Ti},
\mathsf{Hz},\mathsf{Si},\mathsf{Sc},\mathsf{Mo}).
\]
For $r\in\ZZ_7$, set
\[
\mathsf M_r=\{\{\infty,r\}\}\cup
\{\{r+j,r-j\}:j=1,2,3\}.
\]
Then $E(K_8)=\bigsqcup_{r\in\ZZ_7}\mathsf M_r$, and the canonical matching is
\[
\mathsf M_0=\{\mathsf{Gravity\! -\! Light},\;
\mathsf{Memory\! -\! Motion},\;
\mathsf{Time\! -\! Scope},\;
\mathsf{Horizon\! -\! Silence}\}.
\]
}

\hypersetup{pdftitle={Horizon Silence: Recognition at the Boundary of Type},pdfauthor={Anthony Vito Coppa},pdfsubject={Viability horizon, physical occurrence, zero residual, and future difference}}
\begin{document}
\papertitle{Horizon Silence}{Viability boundary and zero residual}

\begin{abstract}
A horizon is the boundary of a physically viable region. Silence is the zero
residual relation at a completed address. A present occurrence has a carrier,
a constraint, a read, a calibration, and a measured difference. Time orders
occurrences into history; future differences remain individually addressable.
The eight-address field is partitioned into seven matchings. Its triangular
residual condition reconstructs an edge potential exactly. The resulting
equations type survival, boundary, observation, completion, and continuation
inside one physical field.
\end{abstract}

\section{Viability, horizon, and silence}

Let $X$ be the state carrier of a living physical system and let
$\mathcal V\subset X$ be its viable region. Its survival condition is
\[
\theta(x)=1\Longleftrightarrow x\in\mathcal V,
\tag{HS.0}
\]
and its horizon is the boundary
\[
\mathcal H_{\rm hor}=\partial\mathcal V.
\tag{HS.1}
\]
For a time-dependent state, regulation carries the system by
\[
\dot x=f(x,u,d),
\qquad \theta(x(t))=1.
\tag{HS.2}
\]
The control $u$ and disturbance $d$ change the trajectory. The viable set
and its boundary determine which changes remain survivable.

\section{Energy and regulation}

For a thermodynamic state with internal energy $U_{\rm th}$, entropy
$S_{\rm th}$, volume $V$, pressure $p$, chemical potentials $\mu_i$, particle
counts $N_i$, and generalized coordinate--force pairs $(X_a,Y_a)$,
\[
dU_{\rm th}=T_{\rm th}\,dS_{\rm th}-p\,dV
+\sum_i\mu_i\,dN_i+\sum_aY_a\,dX_a.
\tag{HS.3}
\]
For an irreversible process,
\[
dS_{\rm th}=d_eS_{\rm th}+d_iS_{\rm th},
\qquad d_iS_{\rm th}\ge0.
\tag{HS.4}
\]
Energy, matter, work, and entropy therefore enter the viable trajectory as
measurable state relations. Regulation changes their flows while preserving
the survival condition in HS.0.

\section{Constraint and type}

Fix a jurisdiction $J$ and its active context
\[
\Gamma_J=(\operatorname{Condition},\operatorname{Constraint},
\operatorname{SourceAuthority},\operatorname{Namespace}).
\tag{HS.5}
\]
The namespace reserves
\[
\varnothing=\mathsf{UNTYPEABLE}.
\tag{HS.6}
\]
A candidate belongs to the typed domain only when it has a type in the active
context:
\[
D_{\Gamma_J}=
\{\sigma:\exists\tau\;(\Gamma_J\vdash\sigma:\tau)\}.
\tag{HS.7}
\]
The gate acts after type:
\[
\Gate_{\Gamma_J}:D_{\Gamma_J}\to\{\mathsf{can},\mathbf2\}.
\tag{HS.8}
\]
Here $\mathbf2$ is the declared fold/refusal value. Set-theoretic emptiness,
when needed, is written $\emptyset_{\rm set}$.

\section{Occurrence, history, and future difference}

\sharedoccurrence{HS.9}{HS.10}

The lower-case $\theta_J(x)$ is the survival condition; the upper-case
$\Theta$ is the occurrence timing/address. The calibration $\kappa$ and
residual $\boldsymbol\varepsilon$ fix the resolution of the read. Put
\[
\mathscr O:=\bigsqcup_J\mathsf{Occ}_J,
\qquad
\mathscr L_n:=\mathscr O^n,
\qquad
\mathcal L_n=(\mathsf{Occ}_1,\ldots,\mathsf{Occ}_n)\in\mathscr L_n.
\tag{HS.11}
\]
Let $\mathcal P_n$ be a candidate subset of $\mathscr O$, and put
\[
\mathcal P_n^{\rm typed}:=
\{p\in\mathcal P_n:\Typed_{\Gamma_{J(p)}}(p)\},
\qquad
\mathbf F_n:=\{\Future(p):p\in\mathcal P_n^{\rm typed}\}.
\tag{HS.12}
\]
The tagged elements of $\mathbf F_n$ remain distinct by their occurrence
addresses. With
$\operatorname{Past}(\mathcal L_n):=\{\mathsf{Occ}_1,\ldots,\mathsf{Occ}_n\}$,
the ordered carrier is
\[
\widehat K_n=\operatorname{Past}(\mathcal L_n)\sqcup\mathbf F_n.
\tag{HS.13}
\]

\begin{definition}[Occurrence]
The next occurrence enters history through the partial map
\[
\Admit_n:\mathbf F_n\rightharpoonup\mathscr L_{n+1},
\tag{HS.14}
\]
For $p\in\mathscr O$, let $x(p)$ denote its state coordinate.
with
\[
\Admit_n(\Future(p))
=\snoc(\mathcal L_n,p)
\Longleftrightarrow
\Gate_{\Gamma_{J(p)}}(p)=\mathsf{can}
\land\mathsf{Complete}(p)
\land\theta_{J(p)}(x(p))=1.
\tag{HS.15}
\]
\end{definition}

The map carries one completed occurrence into history and preserves the order
of every prior occurrence.

\section{Completion and silence}

For $p\in\mathscr O$, define its residual vector by
\[
\boldsymbol\varepsilon(p):=
\bigl(\varepsilon_{\rm drop}(p),
\varepsilon_{\rm add}(p),
\varepsilon_{\rm strength}(p)\bigr)
\in\RR_{\ge0}^{3}.
\tag{HS.16}
\]
For $p=\mathsf{Occ}_{J,n}$, the residual slot in HS.10 is
$\boldsymbol\varepsilon_n=\boldsymbol\varepsilon(p)$.
For an occurrence $p$, the completion condition is
\[
\mathsf{Complete}(p)
\Longleftrightarrow
\Typed_{\Gamma_{J(p)}}(p)
\land
\boldsymbol\varepsilon(p)=0.
\tag{HS.17}
\]
If $\operatorname{Out}_n$ is the physical output set, the local stopping
condition is
\[
\mathsf{Enough}_n
\Longleftrightarrow
\mathsf{Complete}(\mathsf{Occ}_n)
\land \operatorname{Out}_n=\{y_n\}.
\tag{HS.18}
\]

Silence is the local fit reading
\[
\SilFit(p)=1
\Longleftrightarrow
\mathsf{Complete}(p)
\qquad(p\in\mathcal P_n^{\rm typed}).
\tag{HS.19}
\]
The condition states that the required relation is complete at its present
address and reader resolution.

\section{Eight-address field}

\sharedmatching
The canonical matching gives each office one paired relation. The complete
$K_8$ field carries all $28$ unordered pair-relations across seven disjoint
matchings.

\begin{center}
\small
\begin{tabularx}{0.97\linewidth}{@{}lX@{}}
\toprule
$M_0$ & Gr--Li, Me--Mo, Ti--Sc, Hz--Si\\
$M_1$ & Gr--Me, Ti--Li, Hz--Mo, Si--Sc\\
$M_2$ & Gr--Ti, Hz--Me, Si--Li, Sc--Mo\\
$M_3$ & Gr--Hz, Si--Ti, Sc--Me, Mo--Li\\
$M_4$ & Gr--Si, Sc--Hz, Mo--Ti, Li--Me\\
$M_5$ & Gr--Sc, Mo--Si, Li--Hz, Me--Ti\\
$M_6$ & Gr--Mo, Li--Sc, Me--Si, Ti--Hz\\
\bottomrule
\end{tabularx}
\end{center}

Let $d_{ij}\in\RR$ be antisymmetric edge values,
$d_{ji}=-d_{ij}$, and define the triangular residual by
\[
h_{ijk}:=d_{ij}+d_{jk}+d_{ki}.
\tag{HS.20}
\]
Each layer is a four-edge restriction of the complete edge field:
\[
d^{(r)}:=d\vert_{\mathsf M_r},
\qquad
d=(d^{(0)},\ldots,d^{(6)}).
\tag{HS.21}
\]

\begin{theorem}[Potential reconstruction]
For an antisymmetric edge field on $K_8$,
\[
\bigl(\forall\,i,j,k\in\mathcal A_8\text{ distinct}\quad h_{ijk}=0\bigr)
\Longleftrightarrow
\bigl(\forall\,0\le i<j\le6\quad h_{\infty i j}=0\bigr)
\Longleftrightarrow
\bigl(\exists s:\mathcal A_8\to\RR\quad d_{ij}=s_j-s_i\bigr).
\tag{HS.22}
\]
The potential $s$ is unique up to addition of one constant.
\end{theorem}

\begin{proof}
The first condition contains the $21$ triangles through $\infty$. Conversely,
set $s_\infty=0$ and $s_i=d_{\infty i}$. Then
\[
h_{\infty i j}=d_{\infty i}+d_{ij}-d_{\infty j}=0
\]
gives $d_{ij}=s_j-s_i$, and this potential telescopes around every triangle.
Adding a constant to all $s_i$ leaves every difference unchanged; equality of
all edge differences leaves only that freedom.
\end{proof}

Writing $\mathbf h(d)=(h_{ijk})_{\{i,j,k\}\in\binom{\mathcal A_8}{3}}$,
the $21$ residuals through $\infty$ are independent and exhaust the rank:
\[
\operatorname{rank}\mathbf h=21,
\qquad
\dim\ker\mathbf h=28-21=7.
\tag{HS.23}
\]
Each residual $h_{\infty i j}$ contains its own finite edge coordinate
$d_{ij}$, which gives independence. The seven-dimensional kernel is the
potential space modulo its additive constant.

\section{Cyclic matching transport}

The cyclic transport fixes the gravity address and advances every finite
address by one:
\[
\tau(\infty)=\infty,
\qquad
\tau(k)=k+1\pmod 7.
\tag{HS.24}
\]
It acts on an edge by its endpoints. For every $r\in\ZZ_7$,
\[
\tau(\mathsf M_r)=\mathsf M_{r+1},
\qquad
\tau^7=I.
\tag{HS.25}
\]

\begin{proof}
The distinguished edge obeys
\[
\tau\{\infty,r\}=\{\infty,r+1\}.
\]
For $j=1,2,3$,
\[
\tau\{r+j,r-j\}=\{r+1+j,r+1-j\}.
\]
Thus every edge of $\mathsf M_r$ maps to the corresponding edge of
$\mathsf M_{r+1}$. Seven applications return every finite address to itself.
\end{proof}

\section{Cyclic carrier}

\sharedroot

The signed root actions form the order-eight dihedral group of a square:
\[
\langle S,H\rangle
=\{\pm I,\pm S,\pm H,\pm C\}.
\tag{HS.26}
\]
The address carrier $\mathcal A_8$ and the root action carrier are separately
typed. Their common cyclic coordinate is the group element $\tau^r$ that
advances the matching layer $\mathsf M_r$.

Let $E=\RR^2$ carry the Euclidean form $I_2$ and define
\[
\widetilde V=\RR^2\oplus E,
\qquad
\widetilde G=G\oplus I_2,
\qquad
\widetilde C=C\oplus C,
\qquad
\operatorname{sig}(\widetilde G)=(3,1).
\tag{HS.27}
\]
With $D=\diag(1,-1)$, define the extended mirror and cyclic action by
\[
\widetilde H=H\oplus D,
\qquad
R_r=e^{(2\pi r/7)C},
\qquad
\varrho(\tau^r)=I_2\oplus R_r
\quad(r\in\ZZ_7).
\tag{HS.28}
\]
They satisfy
\[
\begin{gathered}
\varrho(\tau^r)\varrho(\tau^s)=\varrho(\tau^{r+s}),
\qquad
\varrho(\tau)^7=I_4,
\qquad
\varrho(\tau)^{\mathsf T}\widetilde G\varrho(\tau)=\widetilde G,\\
\varrho(\tau)^{-1}\widetilde C\varrho(\tau)=\widetilde C,
\qquad
\widetilde H\varrho(\tau)\widetilde H=\varrho(\tau)^{-1}.
\end{gathered}
\tag{HS.29}
\]
Thus $\mathsf M_r\mapsto\varrho(\tau^r)$ carries the common sevenfold layer
coordinate between the finite ledger and the Lorentzian operator carrier.

\end{document}
